
\documentclass{article}
\usepackage{colortbl}
\usepackage{makecell}
\usepackage{multirow}
\usepackage{supertabular}

\begin{document}

\newcounter{utterance}

\centering \large Interaction Transcript for game `matchit\_ascii', experiment `same\_grid', episode 0 with mistral{-}large{-}3{-}azure.
\vspace{24pt}

{ \footnotesize  \setcounter{utterance}{1}
\setlength{\tabcolsep}{0pt}
\begin{supertabular}{c@{$\;$}|p{.15\linewidth}@{}p{.15\linewidth}p{.15\linewidth}p{.15\linewidth}p{.15\linewidth}p{.15\linewidth}}
    \# & $\;$A & \multicolumn{4}{c}{Game Master} & $\;\:$B\\
    \hline

    \theutterance \stepcounter{utterance}  
    & & \multicolumn{4}{p{0.6\linewidth}}{
        \cellcolor[rgb]{0.9,0.9,0.9}{
            \makecell[{{p{\linewidth}}}]{
                \texttt{\tiny{[P1$\langle$GM]}}
                \texttt{Sudjeluješ u igri pogađanja u suradnji.} \\
\texttt{Cilj je otkriti jesu li ova mreža i druga slika koju samo ja mogu vidjeti iste.} \\
\texttt{Molim te, prvo opiši svoju mrežu. Zatim ću ja dati svoj opis i možemo si međusobno postavljati pitanja o mrežama kako bismo utvrdili jesu li iste. Sada započni svoj kratki opis mreže s "OPIS:" nakon čega slijedi opis. Ne dodaj ništa drugo. Ovo je tvoja mreža:} \\
\texttt{▢ X X X X} \\
\texttt{▢ X X X X} \\
\texttt{▢ X X X X} \\
\texttt{▢ X X X X} \\
\texttt{▢ ▢ ▢ ▢ ▢} \\
            }
        }
    }
    & \\ \\

    \theutterance \stepcounter{utterance}  
    & \multicolumn{4}{p{0.6\linewidth}}{
        \cellcolor[rgb]{0.9,0.9,0.9}{
            \makecell[{{p{\linewidth}}}]{
                \texttt{\tiny{[P1$\rangle$GM]}}
                \texttt{OPIS: Mreža je pravokutnik sastavljen od 5 redova i 5 stupaca. U prva četiri retka, prvi kvadrat (s lijeve strane) je prazan (označen s ▢), dok su preostala četiri kvadrata u svakom retku ispunjena s "X". U petom retku svi kvadrati su prazni (▢).} \\
            }
        }
    }
    & & \\ \\

    \theutterance \stepcounter{utterance}  
    & & & \multicolumn{2}{p{0.3\linewidth}}{
        \cellcolor[rgb]{0.9,0.9,0.9}{
            \makecell[{{p{\linewidth}}}]{
                \texttt{\tiny{[GM$|$GM]}}
                \texttt{continue} \\
            }
        }
    }
    & & \\ \\

    \theutterance \stepcounter{utterance}  
    & & \multicolumn{4}{p{0.6\linewidth}}{
        \cellcolor[rgb]{0.9,0.9,0.9}{
            \makecell[{{p{\linewidth}}}]{
                \texttt{\tiny{[GM$\rangle$P2]}}
                \texttt{Sudjeluješ u igri pogađanja u suradnji.} \\
\texttt{Cilj je otkriti jesu li ova mreža i druga slika koju samo ja mogu vidjeti iste.} \\
\texttt{Molim te, prvo opiši svoju mrežu. Zatim ću ja dati svoj opis i možemo si međusobno postavljati pitanja o mrežama kako bismo utvrdili jesu li iste. Sada započni svoj kratki opis mreže s "OPIS:" nakon čega slijedi opis. Ne dodaj ništa drugo. Ovo je tvoja mreža:} \\
\texttt{▢ X X X X} \\
\texttt{▢ X X X X} \\
\texttt{▢ X X X X} \\
\texttt{▢ X X X X} \\
\texttt{▢ ▢ ▢ ▢ ▢} \\
            }
        }
    }
    & \\ \\

    \theutterance \stepcounter{utterance}  
    & & & \multicolumn{4}{p{0.6\linewidth}}{
        \cellcolor[rgb]{0.9,0.9,0.9}{
            \makecell[{{p{\linewidth}}}]{
                \texttt{\tiny{[GM$\langle$P2]}}
                \texttt{OPIS: Mreža je pravokutna, dimenzija 5x5. U prvom stupcu nalaze se prazni kvadratići (▢) u svih pet redova. U preostalim četiri stupca (drugi do peti) nalaze se križići (X) u prva četiri reda, dok je u petom redu svih pet kvadratića prazno (▢).} \\
            }
        }
    }
     \\ \\

    \theutterance \stepcounter{utterance}  
    & & & \multicolumn{2}{p{0.3\linewidth}}{
        \cellcolor[rgb]{0.9,0.9,0.9}{
            \makecell[{{p{\linewidth}}}]{
                \texttt{\tiny{[GM$|$GM]}}
                \texttt{continue} \\
            }
        }
    }
    & & \\ \\

    \theutterance \stepcounter{utterance}  
    & & \multicolumn{4}{p{0.6\linewidth}}{
        \cellcolor[rgb]{0.9,0.9,0.9}{
            \makecell[{{p{\linewidth}}}]{
                \texttt{\tiny{[GM$\rangle$P2]}}
                \texttt{Ovo je mojOPIS: Mreža je pravokutnik sastavljen od 5 redova i 5 stupaca. U prva četiri retka, prvi kvadrat (s lijeve strane) je prazan (označen s ▢), dok su preostala četiri kvadrata u svakom retku ispunjena s "X". U petom retku svi kvadrati su prazni (▢).} \\
\texttt{Sada postavi pitanje kako bi saznao nove aspekte moje mreže koji se možda razlikuju od tvoje mreže. Započni s "PITANJE:" i ne dodavaj ništa drugo.} \\
            }
        }
    }
    & \\ \\

    \theutterance \stepcounter{utterance}  
    & & & \multicolumn{4}{p{0.6\linewidth}}{
        \cellcolor[rgb]{0.9,0.9,0.9}{
            \makecell[{{p{\linewidth}}}]{
                \texttt{\tiny{[GM$\langle$P2]}}
                \texttt{PITANJE: Kako su raspoređeni prazni kvadratići (▢) u petom retku tvoje mreže – jesu li svi poravnati u istom stupcu kao i prazni kvadratići u prva četiri retka, ili postoje odstupanja?} \\
            }
        }
    }
     \\ \\

    \theutterance \stepcounter{utterance}  
    & & & \multicolumn{2}{p{0.3\linewidth}}{
        \cellcolor[rgb]{0.9,0.9,0.9}{
            \makecell[{{p{\linewidth}}}]{
                \texttt{\tiny{[GM$|$GM]}}
                \texttt{continue} \\
            }
        }
    }
    & & \\ \\

    \theutterance \stepcounter{utterance}  
    & & \multicolumn{4}{p{0.6\linewidth}}{
        \cellcolor[rgb]{0.9,0.9,0.9}{
            \makecell[{{p{\linewidth}}}]{
                \texttt{\tiny{[P1$\langle$GM]}}
                \texttt{Ovo je mojOPIS: Mreža je pravokutna, dimenzija 5x5. U prvom stupcu nalaze se prazni kvadratići (▢) u svih pet redova. U preostalim četiri stupca (drugi do peti) nalaze se križići (X) u prva četiri reda, dok je u petom redu svih pet kvadratića prazno (▢).} \\
\texttt{PITANJE: Kako su raspoređeni prazni kvadratići (▢) u petom retku tvoje mreže – jesu li svi poravnati u istom stupcu kao i prazni kvadratići u prva četiri retka, ili postoje odstupanja?Započni svoj odgovor s "ODGOVOR:" i nemoj dodati ništa drugo.} \\
            }
        }
    }
    & \\ \\

    \theutterance \stepcounter{utterance}  
    & \multicolumn{4}{p{0.6\linewidth}}{
        \cellcolor[rgb]{0.9,0.9,0.9}{
            \makecell[{{p{\linewidth}}}]{
                \texttt{\tiny{[P1$\rangle$GM]}}
                \texttt{ODGOVOR: U petom retku moje mreže svi prazni kvadratići (▢) su poravnani u svih pet stupaca, dok su u prva četiri retka prazni samo u prvom stupcu.} \\
            }
        }
    }
    & & \\ \\

    \theutterance \stepcounter{utterance}  
    & & & \multicolumn{2}{p{0.3\linewidth}}{
        \cellcolor[rgb]{0.9,0.9,0.9}{
            \makecell[{{p{\linewidth}}}]{
                \texttt{\tiny{[GM$|$GM]}}
                \texttt{continue} \\
            }
        }
    }
    & & \\ \\

    \theutterance \stepcounter{utterance}  
    & & \multicolumn{4}{p{0.6\linewidth}}{
        \cellcolor[rgb]{0.9,0.9,0.9}{
            \makecell[{{p{\linewidth}}}]{
                \texttt{\tiny{[P1$\langle$GM]}}
                \texttt{Sada postavi pitanje kako bi saznao nove aspekte moje mreže koji se možda razlikuju od tvoje mreže. Započni s "PITANJE:" i ne dodavaj ništa drugo.} \\
            }
        }
    }
    & \\ \\

    \theutterance \stepcounter{utterance}  
    & \multicolumn{4}{p{0.6\linewidth}}{
        \cellcolor[rgb]{0.9,0.9,0.9}{
            \makecell[{{p{\linewidth}}}]{
                \texttt{\tiny{[P1$\rangle$GM]}}
                \texttt{PITANJE: Jesu li u tvojoj mreži prazni kvadratići (▢) u prvom stupcu svih pet redova potpuno poravnani (tj. svi ▢ u istom stupcu), ili postoje redovi gdje se prazni kvadrat pomiče u drugi stupac?} \\
            }
        }
    }
    & & \\ \\

    \theutterance \stepcounter{utterance}  
    & & & \multicolumn{2}{p{0.3\linewidth}}{
        \cellcolor[rgb]{0.9,0.9,0.9}{
            \makecell[{{p{\linewidth}}}]{
                \texttt{\tiny{[GM$|$GM]}}
                \texttt{continue} \\
            }
        }
    }
    & & \\ \\

    \theutterance \stepcounter{utterance}  
    & & \multicolumn{4}{p{0.6\linewidth}}{
        \cellcolor[rgb]{0.9,0.9,0.9}{
            \makecell[{{p{\linewidth}}}]{
                \texttt{\tiny{[GM$\rangle$P2]}}
                \texttt{ODGOVOR: U petom retku moje mreže svi prazni kvadratići (▢) su poravnani u svih pet stupaca, dok su u prva četiri retka prazni samo u prvom stupcu.} \\
\texttt{PITANJE: Jesu li u tvojoj mreži prazni kvadratići (▢) u prvom stupcu svih pet redova potpuno poravnani (tj. svi ▢ u istom stupcu), ili postoje redovi gdje se prazni kvadrat pomiče u drugi stupac?Započni svoj odgovor s "ODGOVOR:" i nemoj dodati ništa drugo.} \\
            }
        }
    }
    & \\ \\

    \theutterance \stepcounter{utterance}  
    & & & \multicolumn{4}{p{0.6\linewidth}}{
        \cellcolor[rgb]{0.9,0.9,0.9}{
            \makecell[{{p{\linewidth}}}]{
                \texttt{\tiny{[GM$\langle$P2]}}
                \texttt{ODGOVOR: Da, u mojoj mreži su prazni kvadratići (▢) u prvom stupcu svih pet redova potpuno poravnani – svi se nalaze u istom stupcu bez pomaka.} \\
            }
        }
    }
     \\ \\

    \theutterance \stepcounter{utterance}  
    & & & \multicolumn{2}{p{0.3\linewidth}}{
        \cellcolor[rgb]{0.9,0.9,0.9}{
            \makecell[{{p{\linewidth}}}]{
                \texttt{\tiny{[GM$|$GM]}}
                \texttt{continue} \\
            }
        }
    }
    & & \\ \\

    \theutterance \stepcounter{utterance}  
    & & \multicolumn{4}{p{0.6\linewidth}}{
        \cellcolor[rgb]{0.9,0.9,0.9}{
            \makecell[{{p{\linewidth}}}]{
                \texttt{\tiny{[GM$\rangle$P2]}}
                \texttt{Sada postavi pitanje kako bi saznao nove aspekte moje mreže koji se možda razlikuju od tvoje mreže. Započni s "PITANJE:" i ne dodavaj ništa drugo.} \\
            }
        }
    }
    & \\ \\

    \theutterance \stepcounter{utterance}  
    & & & \multicolumn{4}{p{0.6\linewidth}}{
        \cellcolor[rgb]{0.9,0.9,0.9}{
            \makecell[{{p{\linewidth}}}]{
                \texttt{\tiny{[GM$\langle$P2]}}
                \texttt{PITANJE: Postoji li u tvojoj mreži neki redak ili stupac gdje se križići (X) i prazni kvadratići (▢) izmjenjuju u uzorku koji se razlikuje od ostalih redaka/stupaca?} \\
            }
        }
    }
     \\ \\

    \theutterance \stepcounter{utterance}  
    & & & \multicolumn{2}{p{0.3\linewidth}}{
        \cellcolor[rgb]{0.9,0.9,0.9}{
            \makecell[{{p{\linewidth}}}]{
                \texttt{\tiny{[GM$|$GM]}}
                \texttt{continue} \\
            }
        }
    }
    & & \\ \\

    \theutterance \stepcounter{utterance}  
    & & \multicolumn{4}{p{0.6\linewidth}}{
        \cellcolor[rgb]{0.9,0.9,0.9}{
            \makecell[{{p{\linewidth}}}]{
                \texttt{\tiny{[P1$\langle$GM]}}
                \texttt{ODGOVOR: Da, u mojoj mreži su prazni kvadratići (▢) u prvom stupcu svih pet redova potpuno poravnani – svi se nalaze u istom stupcu bez pomaka.} \\
\texttt{PITANJE: Postoji li u tvojoj mreži neki redak ili stupac gdje se križići (X) i prazni kvadratići (▢) izmjenjuju u uzorku koji se razlikuje od ostalih redaka/stupaca?Započni svoj odgovor s "ODGOVOR:" i nemoj dodati ništa drugo.} \\
            }
        }
    }
    & \\ \\

    \theutterance \stepcounter{utterance}  
    & \multicolumn{4}{p{0.6\linewidth}}{
        \cellcolor[rgb]{0.9,0.9,0.9}{
            \makecell[{{p{\linewidth}}}]{
                \texttt{\tiny{[P1$\rangle$GM]}}
                \texttt{ODGOVOR: Ne, u mojoj mreži nema retka ni stupca gdje se križići (X) i prazni kvadratići (▢) izmjenjuju u drugom uzorku – svi redovi osim petog slijede isti obrazac (prvi ▢, zatim četiri X), a peti redak je potpuno prazan.} \\
            }
        }
    }
    & & \\ \\

    \theutterance \stepcounter{utterance}  
    & & & \multicolumn{2}{p{0.3\linewidth}}{
        \cellcolor[rgb]{0.9,0.9,0.9}{
            \makecell[{{p{\linewidth}}}]{
                \texttt{\tiny{[GM$|$GM]}}
                \texttt{continue} \\
            }
        }
    }
    & & \\ \\

    \theutterance \stepcounter{utterance}  
    & & \multicolumn{4}{p{0.6\linewidth}}{
        \cellcolor[rgb]{0.9,0.9,0.9}{
            \makecell[{{p{\linewidth}}}]{
                \texttt{\tiny{[P1$\langle$GM]}}
                \texttt{Sada postavi pitanje kako bi saznao nove aspekte moje mreže koji se možda razlikuju od tvoje mreže. Započni s "PITANJE:" i ne dodavaj ništa drugo.} \\
            }
        }
    }
    & \\ \\

    \theutterance \stepcounter{utterance}  
    & \multicolumn{4}{p{0.6\linewidth}}{
        \cellcolor[rgb]{0.9,0.9,0.9}{
            \makecell[{{p{\linewidth}}}]{
                \texttt{\tiny{[P1$\rangle$GM]}}
                \texttt{PITANJE: Da li se u tvojoj mreži u **petom retku** pojavljuje **barem jedan X**, ili je taj redak zaista potpuno prazan (svih pet ▢) kao u mom opisu?} \\
            }
        }
    }
    & & \\ \\

    \theutterance \stepcounter{utterance}  
    & & & \multicolumn{2}{p{0.3\linewidth}}{
        \cellcolor[rgb]{0.9,0.9,0.9}{
            \makecell[{{p{\linewidth}}}]{
                \texttt{\tiny{[GM$|$GM]}}
                \texttt{continue} \\
            }
        }
    }
    & & \\ \\

    \theutterance \stepcounter{utterance}  
    & & \multicolumn{4}{p{0.6\linewidth}}{
        \cellcolor[rgb]{0.9,0.9,0.9}{
            \makecell[{{p{\linewidth}}}]{
                \texttt{\tiny{[GM$\rangle$P2]}}
                \texttt{ODGOVOR: Ne, u mojoj mreži nema retka ni stupca gdje se križići (X) i prazni kvadratići (▢) izmjenjuju u drugom uzorku – svi redovi osim petog slijede isti obrazac (prvi ▢, zatim četiri X), a peti redak je potpuno prazan.} \\
\texttt{PITANJE: Da li se u tvojoj mreži u **petom retku** pojavljuje **barem jedan X**, ili je taj redak zaista potpuno prazan (svih pet ▢) kao u mom opisu?Započni svoj odgovor s "ODGOVOR:" i nemoj dodati ništa drugo.} \\
            }
        }
    }
    & \\ \\

    \theutterance \stepcounter{utterance}  
    & & & \multicolumn{4}{p{0.6\linewidth}}{
        \cellcolor[rgb]{0.9,0.9,0.9}{
            \makecell[{{p{\linewidth}}}]{
                \texttt{\tiny{[GM$\langle$P2]}}
                \texttt{ODGOVOR: U mojoj mreži je peti redak **potpuno prazan** (svih pet ▢), bez ijednog X.} \\
            }
        }
    }
     \\ \\

    \theutterance \stepcounter{utterance}  
    & & & \multicolumn{2}{p{0.3\linewidth}}{
        \cellcolor[rgb]{0.9,0.9,0.9}{
            \makecell[{{p{\linewidth}}}]{
                \texttt{\tiny{[GM$|$GM]}}
                \texttt{continue} \\
            }
        }
    }
    & & \\ \\

    \theutterance \stepcounter{utterance}  
    & & \multicolumn{4}{p{0.6\linewidth}}{
        \cellcolor[rgb]{0.9,0.9,0.9}{
            \makecell[{{p{\linewidth}}}]{
                \texttt{\tiny{[GM$\rangle$P2]}}
                \texttt{Sada postavi pitanje kako bi saznao nove aspekte moje mreže koji se možda razlikuju od tvoje mreže. Započni s "PITANJE:" i ne dodavaj ništa drugo.} \\
            }
        }
    }
    & \\ \\

    \theutterance \stepcounter{utterance}  
    & & & \multicolumn{4}{p{0.6\linewidth}}{
        \cellcolor[rgb]{0.9,0.9,0.9}{
            \makecell[{{p{\linewidth}}}]{
                \texttt{\tiny{[GM$\langle$P2]}}
                \texttt{PITANJE: Postoji li u tvojoj mreži neki **stupac** (osim prvog) u kojem se pojavljuje **više od jednog praznog kvadratića (▢)**, ili su svi prazni ▢ koncentrirani samo u prvom stupcu i petom retku?} \\
            }
        }
    }
     \\ \\

    \theutterance \stepcounter{utterance}  
    & & & \multicolumn{2}{p{0.3\linewidth}}{
        \cellcolor[rgb]{0.9,0.9,0.9}{
            \makecell[{{p{\linewidth}}}]{
                \texttt{\tiny{[GM$|$GM]}}
                \texttt{continue} \\
            }
        }
    }
    & & \\ \\

    \theutterance \stepcounter{utterance}  
    & & \multicolumn{4}{p{0.6\linewidth}}{
        \cellcolor[rgb]{0.9,0.9,0.9}{
            \makecell[{{p{\linewidth}}}]{
                \texttt{\tiny{[P1$\langle$GM]}}
                \texttt{ODGOVOR: U mojoj mreži je peti redak **potpuno prazan** (svih pet ▢), bez ijednog X.} \\
\texttt{PITANJE: Postoji li u tvojoj mreži neki **stupac** (osim prvog) u kojem se pojavljuje **više od jednog praznog kvadratića (▢)**, ili su svi prazni ▢ koncentrirani samo u prvom stupcu i petom retku?Započni svoj odgovor s "ODGOVOR:" i nemoj dodati ništa drugo.} \\
            }
        }
    }
    & \\ \\

    \theutterance \stepcounter{utterance}  
    & \multicolumn{4}{p{0.6\linewidth}}{
        \cellcolor[rgb]{0.9,0.9,0.9}{
            \makecell[{{p{\linewidth}}}]{
                \texttt{\tiny{[P1$\rangle$GM]}}
                \texttt{ODGOVOR: U mojoj mreži su svi prazni kvadratići (▢) koncentrirani **samo u prvom stupcu** (prva četiri retka) i **u cijelom petom retku** – u ostalim stupcima (2–5) nema niti jednog ▢ osim u petom retku.} \\
            }
        }
    }
    & & \\ \\

    \theutterance \stepcounter{utterance}  
    & & & \multicolumn{2}{p{0.3\linewidth}}{
        \cellcolor[rgb]{0.9,0.9,0.9}{
            \makecell[{{p{\linewidth}}}]{
                \texttt{\tiny{[GM$|$GM]}}
                \texttt{continue} \\
            }
        }
    }
    & & \\ \\

    \theutterance \stepcounter{utterance}  
    & & \multicolumn{4}{p{0.6\linewidth}}{
        \cellcolor[rgb]{0.9,0.9,0.9}{
            \makecell[{{p{\linewidth}}}]{
                \texttt{\tiny{[P1$\langle$GM]}}
                \texttt{Sada postavi pitanje kako bi saznao nove aspekte moje mreže koji se možda razlikuju od tvoje mreže. Započni s "PITANJE:" i ne dodavaj ništa drugo.} \\
            }
        }
    }
    & \\ \\

    \theutterance \stepcounter{utterance}  
    & \multicolumn{4}{p{0.6\linewidth}}{
        \cellcolor[rgb]{0.9,0.9,0.9}{
            \makecell[{{p{\linewidth}}}]{
                \texttt{\tiny{[P1$\rangle$GM]}}
                \texttt{PITANJE: Jesu li u tvojoj mreži **svi križići (X)** u stupcima 2–5 **potpuno poravnani** (tj. tvore savršeni pravokutnik bez praznina), ili postoji neki redak ili stupac gdje se X prekida ili pomiče?} \\
            }
        }
    }
    & & \\ \\

    \theutterance \stepcounter{utterance}  
    & & & \multicolumn{2}{p{0.3\linewidth}}{
        \cellcolor[rgb]{0.9,0.9,0.9}{
            \makecell[{{p{\linewidth}}}]{
                \texttt{\tiny{[GM$|$GM]}}
                \texttt{continue} \\
            }
        }
    }
    & & \\ \\

    \theutterance \stepcounter{utterance}  
    & & \multicolumn{4}{p{0.6\linewidth}}{
        \cellcolor[rgb]{0.9,0.9,0.9}{
            \makecell[{{p{\linewidth}}}]{
                \texttt{\tiny{[GM$\rangle$P2]}}
                \texttt{ODGOVOR: U mojoj mreži su svi prazni kvadratići (▢) koncentrirani **samo u prvom stupcu** (prva četiri retka) i **u cijelom petom retku** – u ostalim stupcima (2–5) nema niti jednog ▢ osim u petom retku.} \\
\texttt{PITANJE: Jesu li u tvojoj mreži **svi križići (X)** u stupcima 2–5 **potpuno poravnani** (tj. tvore savršeni pravokutnik bez praznina), ili postoji neki redak ili stupac gdje se X prekida ili pomiče?Započni svoj odgovor s "ODGOVOR:" i nemoj dodati ništa drugo.} \\
            }
        }
    }
    & \\ \\

    \theutterance \stepcounter{utterance}  
    & & & \multicolumn{4}{p{0.6\linewidth}}{
        \cellcolor[rgb]{0.9,0.9,0.9}{
            \makecell[{{p{\linewidth}}}]{
                \texttt{\tiny{[GM$\langle$P2]}}
                \texttt{ODGOVOR: U mojoj mreži su svi križići (X) u stupcima 2–5 **potpuno poravnani** i tvore savršeni pravokutnik od 4×4 (prva četiri retka, stupci 2–5), bez ikakvih prekida ili pomaka. Peti redak je iznimka jer je potpuno prazan.} \\
            }
        }
    }
     \\ \\

    \theutterance \stepcounter{utterance}  
    & & & \multicolumn{2}{p{0.3\linewidth}}{
        \cellcolor[rgb]{0.9,0.9,0.9}{
            \makecell[{{p{\linewidth}}}]{
                \texttt{\tiny{[GM$|$GM]}}
                \texttt{continue} \\
            }
        }
    }
    & & \\ \\

    \theutterance \stepcounter{utterance}  
    & & \multicolumn{4}{p{0.6\linewidth}}{
        \cellcolor[rgb]{0.9,0.9,0.9}{
            \makecell[{{p{\linewidth}}}]{
                \texttt{\tiny{[GM$\rangle$P2]}}
                \texttt{Sada donesi odluku. Što misliš: jesu li tvoja mreža i druga opisana mreža iste? Napiši "ODLUKA: ista mreža" ako misliš da su iste mreže ili "ODLUKA: različite mreže" ako misliš da su različite mreže. Ne dodaj ništa drugo.} \\
            }
        }
    }
    & \\ \\

    \theutterance \stepcounter{utterance}  
    & & & \multicolumn{4}{p{0.6\linewidth}}{
        \cellcolor[rgb]{0.9,0.9,0.9}{
            \makecell[{{p{\linewidth}}}]{
                \texttt{\tiny{[GM$\langle$P2]}}
                \texttt{ODLUKA: ista mreža} \\
            }
        }
    }
     \\ \\

    \theutterance \stepcounter{utterance}  
    & & & \multicolumn{2}{p{0.3\linewidth}}{
        \cellcolor[rgb]{0.9,0.9,0.9}{
            \makecell[{{p{\linewidth}}}]{
                \texttt{\tiny{[GM$|$GM]}}
                \texttt{continue} \\
            }
        }
    }
    & & \\ \\

    \theutterance \stepcounter{utterance}  
    & & & \multicolumn{2}{p{0.3\linewidth}}{
        \cellcolor[rgb]{0.9,0.9,0.9}{
            \makecell[{{p{\linewidth}}}]{
                \texttt{\tiny{[GM$|$GM]}}
                \texttt{success} \\
            }
        }
    }
    & & \\ \\

    \theutterance \stepcounter{utterance}  
    & & \multicolumn{4}{p{0.6\linewidth}}{
        \cellcolor[rgb]{0.9,0.9,0.9}{
            \makecell[{{p{\linewidth}}}]{
                \texttt{\tiny{[P1$\langle$GM]}}
                \texttt{ODGOVOR: U mojoj mreži su svi križići (X) u stupcima 2–5 **potpuno poravnani** i tvore savršeni pravokutnik od 4×4 (prva četiri retka, stupci 2–5), bez ikakvih prekida ili pomaka. Peti redak je iznimka jer je potpuno prazan.} \\
\texttt{Sada donesi odluku. Što misliš: jesu li tvoja mreža i druga opisana mreža iste? Napiši "ODLUKA: ista mreža" ako misliš da su iste mreže ili "ODLUKA: različite mreže" ako misliš da su različite mreže. Ne dodaj ništa drugo.} \\
            }
        }
    }
    & \\ \\

    \theutterance \stepcounter{utterance}  
    & \multicolumn{4}{p{0.6\linewidth}}{
        \cellcolor[rgb]{0.9,0.9,0.9}{
            \makecell[{{p{\linewidth}}}]{
                \texttt{\tiny{[P1$\rangle$GM]}}
                \texttt{ODLUKA: ista mreža} \\
            }
        }
    }
    & & \\ \\

    \theutterance \stepcounter{utterance}  
    & & & \multicolumn{2}{p{0.3\linewidth}}{
        \cellcolor[rgb]{0.9,0.9,0.9}{
            \makecell[{{p{\linewidth}}}]{
                \texttt{\tiny{[GM$|$GM]}}
                \texttt{continue} \\
            }
        }
    }
    & & \\ \\

    \theutterance \stepcounter{utterance}  
    & & & \multicolumn{2}{p{0.3\linewidth}}{
        \cellcolor[rgb]{0.9,0.9,0.9}{
            \makecell[{{p{\linewidth}}}]{
                \texttt{\tiny{[GM$|$GM]}}
                \texttt{success} \\
            }
        }
    }
    & & \\ \\

\end{supertabular}
}

\end{document}
